%%%%%%%%%%%%%%%%%%%%
% Glossary and acronym definitions used by the glossaries package
%%%%%%%%%%%%%%%%%%%%

%%%%%%%%%%
% ACRONYMS
%%%%%%%%%%

%a
\newacronym{abm}{ABM}{Agent Based Modelling}
\newacronym{ace}{ACE}{Agent-based Computational Economics}
\newacronym{ai}{AI}{Artificial Intelligence}

%c
\newacronym{cas}{CAS}{Complex Adaptive Systems}
\newacronym{ces}{CES}{Constant Elasticity of Substitution}

%d
%\newacronym(dsk){DSK}{Dixit Stiglitz Krugman}
\newacronym{dmr}{DMR}{Data Meaning Response}

%e


%g
\newacronym{ge}{GE}{Geographical Economics}

%i
%\newacronym{ied}{IED}{Income Elasticity of Demand}

%l
%\newacronym{lotm}{LOTM}{Law of the Minimum}

%m
%\newacronym{mbf}{MBF}{Movement Buying Factor}

%o
\newacronym{oop}{OOP}{Object-Oriented Programming}

%p
%\newacronym{ped}{PED}{Price Elasticity of Demand}

%%%%%%%%%%%%%%%%%%
% GLOSSARY ENTRIES
%%%%%%%%%%%%%%%%%%

%\newglossaryentry{margrev}{name=marginal revolution,description=
%{``the nearly simultaneous but completely independent discovery in the early 1870s by Jevons, Menger and Walras of the principle of diminishing marginal utility as the fundamental building block of a new kind of static microeconomics.'' (\citealt{blaug_economic_1997} p.277)}}

\newglossaryentry{coremodel}{name=core model,description={Shortened name for Krugman's core-periphery model of regional location.}}

\newglossaryentry{monocomp}{name=monopolistic competition,description={In this thesis, unless otherwise stated, this refers to the approach to imperfect competition developed by Dixit and Stiglitz (1977) in which each firm is guaranteed by their increasing returns to produce one unique variety of good. They have a monopoly on that single good due to increasing returns, but there are enough firms in the market for the price-setting behaviour to not affect others.}}

\newglossaryentry{numeraire}{name=numeraire,description={A product in an economic model that is used as the global measure of value. In the core model, this is agricultural output}}

\newglossaryentry{econgeog}{name=economic geography,description={A broader term than \gls{econgeog}, economic geography refers to any geographical work done on economic problems, whether qualitative or quantitative. In practice, it is now used more often to refer to qualitative approaches}}

\newglossaryentry{geogecon}{name=geographical economics,description={
In this thesis, the term used to refer to work following Paul Krugman's approach. The emphasis is on bringing geography into economics (hence `geographical \emph{economics}', and the methods are those of analytic, neoclassical economics. The term is used by \citeauthor{brakman_new_2009} to replace other variants such as the common 'new economic geography'}}

\newglossaryentry{dxs}{name=Dixit Stiglitz model,description={
Monopolistic Competition model first published by Dixit and Stiglitz in their 1977 paper, 'Monopolistic Competition and Optimum Product Diversity'
}}

\newglossaryentry{repag}{name=representative agent,description={Term used to describe the use of a single agent as a stand-in for all agents in an economic model. The `core model' uses this.}}

\newglossaryentry{bertcomp}{name=Bertrand competition,description={A competitive situation where one firm can take the whole market by lowering their price a tiny amount.}}

\newglossaryentry{lov}{name=love of variety,description={Second of the two forms of `increasing returns'. More formally, elasticity of substitution. High love of variety is low elasticity of substitution: for a given total quantity of goods, consumers prefer a mix of different varieties, and thus are less willing to substitute one for the other.}}

\newglossaryentry{increts}{name=increasing returns,description={Any situation where, for a given set of inputs, the output is more than proportionally larger. Increasing returns are thus non-linear. This thesis uses two types of increasing returns, `love of variety' and `economies of scale' (see other glossary entries).}}

\newglossaryentry{eos}{name=economies of scale,description={First of the two forms of `increasing returns' used in the thesis: a larger input into a production function means a more-than-proportional increase in output.}}


\newglossaryentry{eosub}{name=elasticity of substitution,description={Another way of describing the economic dynamic described by the `Constant Elasticity of Substitution' (CES) function. The other, inversely, is `love of variety'. The more two goods are substitutable for each other, the higher the elasticity of substitution between them, and the lower the `love of variety'. In this thesis, elasticity of substitution only applies to people's utility functions, not inputs into firms' production, so `love of variety' is used.}}





%\newglossaryentry{modelname}{name=SIVA,description = {Java Economic Framework Framework. Got a problem with that?}}
%\newglossaryentry{modelname}{name=SIVA,description = {Name for model framework developed for this thesis: `Spatial Interaction and Visualisation of Agents'. Also, another name for Shiva, destroyer of worlds.}}

%\newglossaryentry{blindsearch}{name=blind search,description={Refers to agents using t\^{a}tonnement or `groping' to optimise - that is, they blindly try some action at $t-1$ in order to find out at $t$ whether it improved a given desired quantity. A catch-all term used to cover all forms of blind optimisation from hill-climbing to full genetic algorithms}}
